\section{Introduction}
\label{sec:intro}

Convolutional neural networks trained on natural image datasets develop
systematic biases toward low-level texture features over global
shape~\citep{geirhos2019texture}.  On fine-grained datasets the effect is
pronounced: classifiers readily latch onto contextual cues such as the grass
behind a terrier, the greenhouse behind a flower, or the sky behind an
aircraft.  These correlations hold in the training distribution and can yield
high validation accuracy, yet they transfer poorly to images where the shortcut
is absent and leave the model brittle under distribution shift.

Data augmentation is the standard remedy.  The prevailing approach applies a
fixed family of photometric and geometric transformations uniformly throughout
training.  AutoAugment~\citep{cubuk2019autoaugment} and its descendants,
RandAugment~\citep{cubuk2020randaugment} and
TrivialAugment~\citep{muller2021trivialaugment}, search over large
transformation spaces but treat the network as a black box: the policy is fixed
before training and never consults the model's internal state.  Spatial-masking
regularisers such as Cutout~\citep{devries2017cutout} and Random
Erasing~\citep{zhong2020randomerasing} remove rectangular regions to discourage
reliance on any single patch, but the location of the mask is chosen uniformly
at random.  A random mask is blind to content; it deletes empty background as
often as it deletes the object.

This raises a natural question.  If a saliency method can tell us which regions
the model currently treats as discriminative, can we place the mask there
deliberately?  KeepAugment~\citep{gong2021keepaugment} explores one side of
this idea by \emph{protecting} high-saliency regions from augmentation, on the
premise that preserving informative content helps accuracy.  We examine the
complementary operations.  What happens when the high-saliency regions are
\emph{removed}, forcing the classifier to find evidence elsewhere?  And what
happens when only the low-saliency regions, the background and irrelevant
texture, are removed, concentrating the perturbation away from the object?

We formalise both operations as augmentation methods.  \textbf{Intelligent
Coarse Dropout (ICD)} masks the highest-saliency tiles of each input;
\textbf{Anti-ICD (AICD)} masks the lowest-saliency tiles.  Both derive their
masks from a GradCAM~\citep{selvaraju2017gradcam} saliency map evaluated at the
model's current parameters, so the augmentation tracks what the model has
learned at each stage of training.  Crucially, the masking decision is made at
the level of coarse tiles rather than individual pixels: the image is divided
into an $8\times 8$-pixel grid (by default), each tile is scored by its mean
saliency, and tiles are masked by a percentile threshold.  Coarse tiling makes
the mask robust to the high-frequency noise that pixel-level saliency maps
often exhibit, while keeping the masked region spatially coherent.  Rather than
a hard black rectangle, masked tiles are filled with a soft strategy: a blurred
copy of the original, a local or global mean colour, matched noise, or a
constant.  The surrounding context then degrades gracefully.

The contributions of this paper are as follows.
\begin{enumerate}
  \item We introduce ICD and AICD, two saliency-conditioned coarse-dropout
    augmentations, and give the formal tile-based mask construction together
    with the percentile-threshold rule that distinguishes them.
  \item We describe five fill strategies for masked tiles and characterise how
    each alters the information that remains available to the model.
  \item We discuss the methods' expected behaviour in terms of curriculum
    learning and spurious-correlation removal, and situate them with respect to
    random masking and saliency-protecting augmentation.
  \item We document a reference implementation in the open-source BNNR
    library~\citep{walo2026bnnr}, including the use of a precomputed saliency
    cache that makes online application practical.
\end{enumerate}

This paper is concerned with the methods themselves.  A controlled empirical
comparison of ICD and AICD against RandAugment, TrivialAugment, and AutoAugment,
together with the accompanying benchmark protocol, is the subject of a separate
study.
